% =============================================================
%  Chapter -- transformer connections in Germany
%  This file holds the chapter heading and its introduction only; the
%  content is one file per subsection, \input below in reading order:
%    vector-groups.tex         IEC vector group and clock notation
%    network-transformers.tex  MV/LV transformers Dyn5, Yzn5, Dyn11
%    lv-lv.tex                 690 V, adapter for US equipment, IT isolating
%    earthing-systems.tex      TN-S, TN-C, TN-C-S, TT, IT
%    comparison.tex            North American and German practice
%  Figures: fig/de_*.tex, written by
%  north-american-practice/fig/connection_figures.py
% =============================================================
\section{Transformer Connections in Germany}
\label{sec:de}

The German low-voltage system is 400/230\,V at 50\,Hz~[12]: a
three-phase four-wire system with 400\,V between lines and 230\,V from line
to neutral. There is no split-phase supply; single-phase loads are
connected between one line and the neutral. The low-voltage network is fed
directly from MV/LV substations at 10 or 20\,kV, and a transformer between
two low-voltage levels is the exception. Customer installations are
connected under VDE-AR-N~4100 at low voltage~[13] and VDE-AR-N~4110
at medium voltage~[14].

Section~\ref{sec:de-vector} gives the vector group notation,
Section~\ref{sec:de-network} the network transformers and
Section~\ref{sec:de-lvlv} the transformers between two low-voltage levels.
Section~\ref{sec:de-earthing} describes the earthing systems, and
Section~\ref{sec:de-comparison} compares German and North American
practice.

\subsection{Vector group notation}
\label{sec:de-vector}

IEC~60076-1 designates a three-phase connection by its vector
group~[2]. A capital letter gives the connection of the
higher-voltage winding (D delta, Y wye, Z zigzag), a lower-case letter that
of the lower-voltage winding, and N or n a neutral that is brought out. The
clock number $h$ gives the phase displacement by which the lower-voltage
side lags the higher-voltage side:
\begin{equation}
  \varphi = h \cdot 30\,^{\circ}
  \label{eq:de-clock}
\end{equation}
The name comes from the clock face: with the higher-voltage phasor of
phase~U at 12, the lower-voltage phasor of the same phase points at the
hour $h$. Dyn5 lags by 150\,$^{\circ}$, Dyn11 by 330\,$^{\circ}$, which is
a lead of 30\,$^{\circ}$, and the North American delta--wye connection with
the lower-voltage side lagging by 30\,$^{\circ}$ is Dyn1.
\FloatBarrier

\subsection{MV/LV network transformers}
\label{sec:de-network}

Figure~\ref{fig:de-network} shows the vector groups of the network
transformer. The clock between the two diagrams gives the vector-group hour
of Equation~\eqref{eq:de-clock}.

\begin{figure}[!ht]
  \centering
  % generated by north-american-practice/fig/connection_figures.py -- do not edit
\begin{tikzpicture}[x=0.27mm, y=-0.27mm, line cap=round, line join=round]
  \draw[hv] (90.0,40.0) -- (142.0,130.0);
  \draw[hv] (142.0,130.0) -- (38.0,130.0);
  \node[font=\footnotesize, text=atInfo, anchor=center, inner sep=0.5pt] at (90.0,145.0) {20\,kV};
  \draw[hv] (38.0,130.0) -- (90.0,40.0);
  \fill[atInfo] (90.0,40.0) circle (3.0);
  \node[font=\footnotesize, text=atInfo, anchor=center, inner sep=0.5pt] at (90.0,27.0) {1U};
  \fill[atInfo] (142.0,130.0) circle (3.0);
  \node[font=\footnotesize, text=atInfo, anchor=center, inner sep=0.5pt] at (153.2,136.5) {1V};
  \fill[atInfo] (38.0,130.0) circle (3.0);
  \node[font=\footnotesize, text=atInfo, anchor=center, inner sep=0.5pt] at (26.8,136.5) {1W};
  \draw[lv] (318.0,100.0) -- (348.0,152.0);
  \draw[lv] (318.0,100.0) -- (258.0,100.0);
  \draw[lv] (318.0,100.0) -- (348.0,48.0);
  \node[font=\footnotesize, text=atAccentFg, anchor=center, inner sep=0.5pt] at (344.3,119.5) {230};
  \fill[atFg] (318.0,100.0) circle (3.0);
  \draw[gnd] (318.0,103.0) -- (318.0,113.0);
  \draw[gnd] (311.0,113.0) -- (325.0,113.0);
  \draw[gnd] (313.5,117.0) -- (322.5,117.0);
  \draw[gnd] (316.0,121.0) -- (320.0,121.0);
  \node[font=\footnotesize, text=atFg, anchor=east, inner sep=0.5pt] at (306.0,120.0) {2N};
  \fill[atAccentFg] (348.0,152.0) circle (3.0);
  \node[font=\footnotesize, text=atAccentFg, anchor=center, inner sep=0.5pt] at (354.5,163.2) {2U};
  \fill[atAccentFg] (258.0,100.0) circle (3.0);
  \node[font=\footnotesize, text=atAccentFg, anchor=center, inner sep=0.5pt] at (245.0,100.0) {2V};
  \fill[atAccentFg] (348.0,48.0) circle (3.0);
  \node[font=\footnotesize, text=atAccentFg, anchor=center, inner sep=0.5pt] at (354.5,36.8) {2W};
  \draw[aux, solid, fill=white] (188.0,58.0) circle (15);
  \draw[aux, solid] (188.0,43.0) -- (188.0,46.0);
  \draw[aux, solid] (195.5,45.0) -- (194.0,47.6);
  \draw[aux, solid] (201.0,50.5) -- (198.4,52.0);
  \draw[aux, solid] (203.0,58.0) -- (200.0,58.0);
  \draw[aux, solid] (201.0,65.5) -- (198.4,64.0);
  \draw[aux, solid] (195.5,71.0) -- (194.0,68.4);
  \draw[aux, solid] (188.0,73.0) -- (188.0,70.0);
  \draw[aux, solid] (180.5,71.0) -- (182.0,68.4);
  \draw[aux, solid] (175.0,65.5) -- (177.6,64.0);
  \draw[aux, solid] (173.0,58.0) -- (176.0,58.0);
  \draw[aux, solid] (175.0,50.5) -- (177.6,52.0);
  \draw[aux, solid] (180.5,45.0) -- (182.0,47.6);
  \draw[hv, line width=0.9pt] (188.0,58.0) -- (188.0,46.0);
  \draw[lv, line width=0.9pt] (188.0,58.0) -- (192.8,66.2);
  \draw[arr] (172.0,104.0) -- (204.0,104.0);
  \node[font=\footnotesize, text=atFgMuted, anchor=center, inner sep=0.5pt] at (188.0,122.0) {150\textdegree};
  \node[font=\footnotesize, text=atFgMuted, anchor=center, inner sep=0.5pt] at (90.0,196.0) {20\,kV $\Delta$};
  \node[font=\footnotesize, text=atFgMuted, anchor=center, inner sep=0.5pt] at (318.0,196.0) {400/230\,V, 2N earthed};
  \node[font=\small\bfseries, text=atFg, anchor=west, inner sep=0.5pt] at (-6.0,18.0) {(a) Dyn5};
  \draw[hv] (90.0,315.0) -- (90.0,255.0);
  \draw[hv] (90.0,315.0) -- (142.0,345.0);
  \draw[hv] (90.0,315.0) -- (38.0,345.0);
  \fill[atInfo] (90.0,255.0) circle (3.0);
  \node[font=\footnotesize, text=atInfo, anchor=center, inner sep=0.5pt] at (90.0,242.0) {1U};
  \fill[atInfo] (142.0,345.0) circle (3.0);
  \node[font=\footnotesize, text=atInfo, anchor=center, inner sep=0.5pt] at (153.2,351.5) {1V};
  \fill[atInfo] (38.0,345.0) circle (3.0);
  \node[font=\footnotesize, text=atInfo, anchor=center, inner sep=0.5pt] at (26.8,351.5) {1W};
  \draw[lv] (318.0,315.0) -- (318.0,349.6) -- (348.0,367.0);
  \draw[lv] (318.0,315.0) -- (288.0,297.7) -- (258.0,315.0);
  \draw[lv] (318.0,315.0) -- (348.0,297.7) -- (348.0,263.0);
  \fill[atFg] (318.0,315.0) circle (3.0);
  \draw[gnd] (318.0,318.0) -- (318.0,328.0);
  \draw[gnd] (311.0,328.0) -- (325.0,328.0);
  \draw[gnd] (313.5,332.0) -- (322.5,332.0);
  \draw[gnd] (316.0,336.0) -- (320.0,336.0);
  \node[font=\footnotesize, text=atFg, anchor=west, inner sep=0.5pt] at (330.0,335.0) {2N};
  \fill[atAccentFg] (348.0,367.0) circle (3.0);
  \node[font=\footnotesize, text=atAccentFg, anchor=center, inner sep=0.5pt] at (354.5,378.2) {2U};
  \fill[atAccentFg] (258.0,315.0) circle (3.0);
  \node[font=\footnotesize, text=atAccentFg, anchor=center, inner sep=0.5pt] at (245.0,315.0) {2V};
  \fill[atAccentFg] (348.0,263.0) circle (3.0);
  \node[font=\footnotesize, text=atAccentFg, anchor=center, inner sep=0.5pt] at (354.5,251.8) {2W};
  \draw[aux, solid, fill=white] (188.0,273.0) circle (15);
  \draw[aux, solid] (188.0,258.0) -- (188.0,261.0);
  \draw[aux, solid] (195.5,260.0) -- (194.0,262.6);
  \draw[aux, solid] (201.0,265.5) -- (198.4,267.0);
  \draw[aux, solid] (203.0,273.0) -- (200.0,273.0);
  \draw[aux, solid] (201.0,280.5) -- (198.4,279.0);
  \draw[aux, solid] (195.5,286.0) -- (194.0,283.4);
  \draw[aux, solid] (188.0,288.0) -- (188.0,285.0);
  \draw[aux, solid] (180.5,286.0) -- (182.0,283.4);
  \draw[aux, solid] (175.0,280.5) -- (177.6,279.0);
  \draw[aux, solid] (173.0,273.0) -- (176.0,273.0);
  \draw[aux, solid] (175.0,265.5) -- (177.6,267.0);
  \draw[aux, solid] (180.5,260.0) -- (182.0,262.6);
  \draw[hv, line width=0.9pt] (188.0,273.0) -- (188.0,261.0);
  \draw[lv, line width=0.9pt] (188.0,273.0) -- (192.8,281.2);
  \draw[arr] (172.0,319.0) -- (204.0,319.0);
  \node[font=\footnotesize, text=atFgMuted, anchor=center, inner sep=0.5pt] at (188.0,337.0) {150\textdegree};
  \node[font=\footnotesize, text=atFgMuted, anchor=center, inner sep=0.5pt] at (90.0,411.0) {20\,kV Y, no neutral};
  \node[font=\footnotesize, text=atFgMuted, anchor=center, inner sep=0.5pt] at (318.0,411.0) {400/230\,V zigzag};
  \node[font=\small\bfseries, text=atFg, anchor=west, inner sep=0.5pt] at (-6.0,233.0) {(b) Yzn5};
  \draw[hv] (90.0,470.0) -- (142.0,560.0);
  \draw[hv] (142.0,560.0) -- (38.0,560.0);
  \node[font=\footnotesize, text=atInfo, anchor=center, inner sep=0.5pt] at (90.0,575.0) {20\,kV};
  \draw[hv] (38.0,560.0) -- (90.0,470.0);
  \fill[atInfo] (90.0,470.0) circle (3.0);
  \node[font=\footnotesize, text=atInfo, anchor=center, inner sep=0.5pt] at (90.0,457.0) {1U};
  \fill[atInfo] (142.0,560.0) circle (3.0);
  \node[font=\footnotesize, text=atInfo, anchor=center, inner sep=0.5pt] at (153.2,566.5) {1V};
  \fill[atInfo] (38.0,560.0) circle (3.0);
  \node[font=\footnotesize, text=atInfo, anchor=center, inner sep=0.5pt] at (26.8,566.5) {1W};
  \draw[lv] (318.0,530.0) -- (288.0,478.0);
  \draw[lv] (318.0,530.0) -- (378.0,530.0);
  \draw[lv] (318.0,530.0) -- (288.0,582.0);
  \node[font=\footnotesize, text=atAccentFg, anchor=center, inner sep=0.5pt] at (291.7,510.5) {230};
  \fill[atFg] (318.0,530.0) circle (3.0);
  \draw[gnd] (318.0,533.0) -- (318.0,543.0);
  \draw[gnd] (311.0,543.0) -- (325.0,543.0);
  \draw[gnd] (313.5,547.0) -- (322.5,547.0);
  \draw[gnd] (316.0,551.0) -- (320.0,551.0);
  \node[font=\footnotesize, text=atFg, anchor=west, inner sep=0.5pt] at (330.0,550.0) {2N};
  \fill[atAccentFg] (288.0,478.0) circle (3.0);
  \node[font=\footnotesize, text=atAccentFg, anchor=center, inner sep=0.5pt] at (281.5,466.8) {2U};
  \fill[atAccentFg] (378.0,530.0) circle (3.0);
  \node[font=\footnotesize, text=atAccentFg, anchor=center, inner sep=0.5pt] at (391.0,530.0) {2V};
  \fill[atAccentFg] (288.0,582.0) circle (3.0);
  \node[font=\footnotesize, text=atAccentFg, anchor=center, inner sep=0.5pt] at (281.5,593.2) {2W};
  \draw[aux, solid, fill=white] (188.0,488.0) circle (15);
  \draw[aux, solid] (188.0,473.0) -- (188.0,476.0);
  \draw[aux, solid] (195.5,475.0) -- (194.0,477.6);
  \draw[aux, solid] (201.0,480.5) -- (198.4,482.0);
  \draw[aux, solid] (203.0,488.0) -- (200.0,488.0);
  \draw[aux, solid] (201.0,495.5) -- (198.4,494.0);
  \draw[aux, solid] (195.5,501.0) -- (194.0,498.4);
  \draw[aux, solid] (188.0,503.0) -- (188.0,500.0);
  \draw[aux, solid] (180.5,501.0) -- (182.0,498.4);
  \draw[aux, solid] (175.0,495.5) -- (177.6,494.0);
  \draw[aux, solid] (173.0,488.0) -- (176.0,488.0);
  \draw[aux, solid] (175.0,480.5) -- (177.6,482.0);
  \draw[aux, solid] (180.5,475.0) -- (182.0,477.6);
  \draw[hv, line width=0.9pt] (188.0,488.0) -- (188.0,476.0);
  \draw[lv, line width=0.9pt] (188.0,488.0) -- (183.2,479.8);
  \draw[arr] (172.0,534.0) -- (204.0,534.0);
  \node[font=\footnotesize, text=atFgMuted, anchor=center, inner sep=0.5pt] at (188.0,552.0) {330\textdegree};
  \node[font=\footnotesize, text=atFgMuted, anchor=center, inner sep=0.5pt] at (90.0,626.0) {20\,kV $\Delta$};
  \node[font=\footnotesize, text=atFgMuted, anchor=center, inner sep=0.5pt] at (318.0,626.0) {400/230\,V, 2N earthed};
  \node[font=\small\bfseries, text=atFg, anchor=west, inner sep=0.5pt] at (-6.0,448.0) {(c) Dyn11};
\end{tikzpicture}

  \caption{Vector groups of MV/LV network transformers 20\,kV --
  400/230\,V: (a) Dyn5, (b) Yzn5, (c) Dyn11. Clock: higher-voltage phasor
  of phase U at 12, lower-voltage phasor at the clock hour.}
  \label{fig:de-network}
\end{figure}

\subsubsection{Dyn5}
Dyn5 is the vector group of German MV/LV network transformers,
Figure~\ref{fig:de-network}(a). The wye secondary is solidly earthed at 2N
and supplies 230\,V from line to neutral. The delta primary provides a
closed path for the zero-sequence current of unbalanced load and of earth
faults on the low-voltage side, so that the zero-sequence impedance seen
from the low-voltage side is low and no zero-sequence current flows in the
medium-voltage network. Network transformers up to 36\,kV are specified to
EN~50588-1~[15], and their losses are limited by the Ecodesign
regulation~[16].

\subsubsection{Yzn5}
In the zigzag secondary of Figure~\ref{fig:de-network}(b), each phase
consists of two half-windings on two different limbs. A single-phase load
then magnetises two limbs, and its ampere-turns are balanced on the primary
side without a delta winding. Yzn5 is used for small network transformers
that supply predominantly single-phase load.

\subsubsection{Dyn11}
Dyn11 has the same windings as Dyn5, connected for a lower-voltage side
leading by 30\,$^{\circ}$, Figure~\ref{fig:de-network}(c). It is the usual
vector group in a number of other countries and occurs in German industrial
installations. A Dyn5 and a Dyn11 transformer are operated in parallel only
after two phases are interchanged on both sides of one of them~[17].
\FloatBarrier

\subsection{Transformers between two low-voltage levels}
\label{sec:de-lvlv}

\subsubsection{690\,V to 400/230\,V}
Industrial installations with large motors distribute at 690\,V, the
line-to-line voltage of the 400/690\,V system of DIN~EN~60038~[12]:
$400\,\mathrm{V} \cdot \sqrt{3} = 692.8\,\mathrm{V}$. A Dyn5 transformer
derives the 400/230\,V supply for auxiliaries and lighting,
Figure~\ref{fig:de-lvlv}(a). The same arrangement occurs in wind turbines
and drive systems.

\begin{figure}[!ht]
  \centering
  % generated by north-american-practice/fig/connection_figures.py -- do not edit
\begin{tikzpicture}[x=0.27mm, y=-0.27mm, line cap=round, line join=round]
  \draw[hv] (90.0,40.0) -- (142.0,130.0);
  \draw[hv] (142.0,130.0) -- (38.0,130.0);
  \node[font=\footnotesize, text=atInfo, anchor=center, inner sep=0.5pt] at (90.0,145.0) {690\,V};
  \draw[hv] (38.0,130.0) -- (90.0,40.0);
  \fill[atInfo] (90.0,40.0) circle (3.0);
  \node[font=\footnotesize, text=atInfo, anchor=center, inner sep=0.5pt] at (90.0,27.0) {1U};
  \fill[atInfo] (142.0,130.0) circle (3.0);
  \node[font=\footnotesize, text=atInfo, anchor=center, inner sep=0.5pt] at (153.2,136.5) {1V};
  \fill[atInfo] (38.0,130.0) circle (3.0);
  \node[font=\footnotesize, text=atInfo, anchor=center, inner sep=0.5pt] at (26.8,136.5) {1W};
  \draw[lv] (318.0,100.0) -- (348.0,152.0);
  \draw[lv] (318.0,100.0) -- (258.0,100.0);
  \draw[lv] (318.0,100.0) -- (348.0,48.0);
  \node[font=\footnotesize, text=atAccentFg, anchor=center, inner sep=0.5pt] at (344.3,119.5) {230};
  \fill[atFg] (318.0,100.0) circle (3.0);
  \draw[gnd] (318.0,103.0) -- (318.0,113.0);
  \draw[gnd] (311.0,113.0) -- (325.0,113.0);
  \draw[gnd] (313.5,117.0) -- (322.5,117.0);
  \draw[gnd] (316.0,121.0) -- (320.0,121.0);
  \node[font=\footnotesize, text=atFg, anchor=east, inner sep=0.5pt] at (306.0,120.0) {2N};
  \fill[atAccentFg] (348.0,152.0) circle (3.0);
  \node[font=\footnotesize, text=atAccentFg, anchor=center, inner sep=0.5pt] at (354.5,163.2) {2U};
  \fill[atAccentFg] (258.0,100.0) circle (3.0);
  \node[font=\footnotesize, text=atAccentFg, anchor=center, inner sep=0.5pt] at (245.0,100.0) {2V};
  \fill[atAccentFg] (348.0,48.0) circle (3.0);
  \node[font=\footnotesize, text=atAccentFg, anchor=center, inner sep=0.5pt] at (354.5,36.8) {2W};
  \draw[aux, solid, fill=white] (188.0,58.0) circle (15);
  \draw[aux, solid] (188.0,43.0) -- (188.0,46.0);
  \draw[aux, solid] (195.5,45.0) -- (194.0,47.6);
  \draw[aux, solid] (201.0,50.5) -- (198.4,52.0);
  \draw[aux, solid] (203.0,58.0) -- (200.0,58.0);
  \draw[aux, solid] (201.0,65.5) -- (198.4,64.0);
  \draw[aux, solid] (195.5,71.0) -- (194.0,68.4);
  \draw[aux, solid] (188.0,73.0) -- (188.0,70.0);
  \draw[aux, solid] (180.5,71.0) -- (182.0,68.4);
  \draw[aux, solid] (175.0,65.5) -- (177.6,64.0);
  \draw[aux, solid] (173.0,58.0) -- (176.0,58.0);
  \draw[aux, solid] (175.0,50.5) -- (177.6,52.0);
  \draw[aux, solid] (180.5,45.0) -- (182.0,47.6);
  \draw[hv, line width=0.9pt] (188.0,58.0) -- (188.0,46.0);
  \draw[lv, line width=0.9pt] (188.0,58.0) -- (192.8,66.2);
  \draw[arr] (172.0,104.0) -- (204.0,104.0);
  \node[font=\footnotesize, text=atFgMuted, anchor=center, inner sep=0.5pt] at (188.0,122.0) {150\textdegree};
  \node[font=\footnotesize, text=atFgMuted, anchor=center, inner sep=0.5pt] at (90.0,196.0) {690\,V $\Delta$};
  \node[font=\footnotesize, text=atFgMuted, anchor=center, inner sep=0.5pt] at (318.0,196.0) {400/230\,V};
  \node[font=\small\bfseries, text=atFg, anchor=west, inner sep=0.5pt] at (-6.0,18.0) {(a)};
  \draw[hv] (90.0,255.0) -- (142.0,345.0);
  \draw[hv] (142.0,345.0) -- (38.0,345.0);
  \node[font=\footnotesize, text=atInfo, anchor=center, inner sep=0.5pt] at (90.0,360.0) {400\,V};
  \draw[hv] (38.0,345.0) -- (90.0,255.0);
  \fill[atInfo] (90.0,255.0) circle (3.0);
  \node[font=\footnotesize, text=atInfo, anchor=center, inner sep=0.5pt] at (90.0,242.0) {1U};
  \fill[atInfo] (142.0,345.0) circle (3.0);
  \node[font=\footnotesize, text=atInfo, anchor=center, inner sep=0.5pt] at (153.2,351.5) {1V};
  \fill[atInfo] (38.0,345.0) circle (3.0);
  \node[font=\footnotesize, text=atInfo, anchor=center, inner sep=0.5pt] at (26.8,351.5) {1W};
  \draw[lv] (318.0,315.0) -- (348.0,367.0);
  \draw[lv] (318.0,315.0) -- (258.0,315.0);
  \draw[lv] (318.0,315.0) -- (348.0,263.0);
  \node[font=\footnotesize, text=atAccentFg, anchor=center, inner sep=0.5pt] at (344.3,334.5) {277};
  \fill[atFg] (318.0,315.0) circle (3.0);
  \draw[gnd] (318.0,318.0) -- (318.0,328.0);
  \draw[gnd] (311.0,328.0) -- (325.0,328.0);
  \draw[gnd] (313.5,332.0) -- (322.5,332.0);
  \draw[gnd] (316.0,336.0) -- (320.0,336.0);
  \node[font=\footnotesize, text=atFg, anchor=east, inner sep=0.5pt] at (306.0,335.0) {2N};
  \fill[atAccentFg] (348.0,367.0) circle (3.0);
  \node[font=\footnotesize, text=atAccentFg, anchor=center, inner sep=0.5pt] at (354.5,378.2) {2U};
  \fill[atAccentFg] (258.0,315.0) circle (3.0);
  \node[font=\footnotesize, text=atAccentFg, anchor=center, inner sep=0.5pt] at (245.0,315.0) {2V};
  \fill[atAccentFg] (348.0,263.0) circle (3.0);
  \node[font=\footnotesize, text=atAccentFg, anchor=center, inner sep=0.5pt] at (354.5,251.8) {2W};
  \draw[aux, solid, fill=white] (188.0,273.0) circle (15);
  \draw[aux, solid] (188.0,258.0) -- (188.0,261.0);
  \draw[aux, solid] (195.5,260.0) -- (194.0,262.6);
  \draw[aux, solid] (201.0,265.5) -- (198.4,267.0);
  \draw[aux, solid] (203.0,273.0) -- (200.0,273.0);
  \draw[aux, solid] (201.0,280.5) -- (198.4,279.0);
  \draw[aux, solid] (195.5,286.0) -- (194.0,283.4);
  \draw[aux, solid] (188.0,288.0) -- (188.0,285.0);
  \draw[aux, solid] (180.5,286.0) -- (182.0,283.4);
  \draw[aux, solid] (175.0,280.5) -- (177.6,279.0);
  \draw[aux, solid] (173.0,273.0) -- (176.0,273.0);
  \draw[aux, solid] (175.0,265.5) -- (177.6,267.0);
  \draw[aux, solid] (180.5,260.0) -- (182.0,262.6);
  \draw[hv, line width=0.9pt] (188.0,273.0) -- (188.0,261.0);
  \draw[lv, line width=0.9pt] (188.0,273.0) -- (192.8,281.2);
  \draw[arr] (172.0,319.0) -- (204.0,319.0);
  \node[font=\footnotesize, text=atFgMuted, anchor=center, inner sep=0.5pt] at (188.0,337.0) {150\textdegree};
  \node[font=\footnotesize, text=atFgMuted, anchor=center, inner sep=0.5pt] at (90.0,411.0) {400\,V $\Delta$};
  \node[font=\footnotesize, text=atFgMuted, anchor=center, inner sep=0.5pt] at (318.0,411.0) {480Y/277\,V};
  \node[font=\small\bfseries, text=atFg, anchor=west, inner sep=0.5pt] at (-6.0,233.0) {(b)};
\end{tikzpicture}

  \caption{Transformers between two low-voltage levels, vector group Dyn5:
  (a) 690\,V -- 400/230\,V, (b) 400\,V -- 480Y/277\,V for equipment built
  to North American ratings.}
  \label{fig:de-lvlv}
\end{figure}

\subsubsection{Supply of equipment built to North American ratings}
A 400\,V -- 480Y/277\,V transformer, Figure~\ref{fig:de-lvlv}(b), or one
with a 208Y/120\,V or 240\,V delta secondary, adapts the voltage; the
frequency remains 50\,Hz. Motors run at $50/60 = 0.833$ of their 60\,Hz
speed. Magnetic equipment rated 480\,V at 60\,Hz operates at
$480\,\mathrm{V} / 50\,\mathrm{Hz} = 9.6\,\mathrm{V/Hz}$ instead of
$8.0\,\mathrm{V/Hz}$, a flux 20\,\% above rated. The voltage at 50\,Hz that
keeps the rated flux is
\begin{equation}
  U_{50} = U_{60} \cdot \frac{50\,\mathrm{Hz}}{60\,\mathrm{Hz}}
         = 480\,\mathrm{V} \cdot \frac{50}{60} = 400\,\mathrm{V}
  \label{eq:de-vhz}
\end{equation}
so that motors rated 480\,V at 60\,Hz are connected to the 400\,V system
directly when the reduced speed is acceptable. The 480\,V transformer
applies to equipment rated for 480\,V at 50\,Hz; equipment that requires
60\,Hz is supplied through a frequency converter.

\subsubsection{Isolating transformer for an IT system}
Group~2 medical locations, such as operating theatres and intensive care,
are supplied through a single-phase isolating transformer whose secondary
is not earthed, Figure~\ref{fig:de-it}~[18]. An insulation
monitoring device to IEC~61557-8~[19] measures the insulation
resistance of the secondary to earth. A first insulation fault raises an
alarm and does not interrupt the supply. DIN~VDE~0100-710 limits the rated
output of the transformer to 0.5 to 10\,kVA~[18].

\begin{figure}[!ht]
  \centering
  % generated by north-american-practice/fig/connection_figures.py -- do not edit
\begin{tikzpicture}[x=0.3mm, y=-0.3mm, line cap=round, line join=round]
  \draw[hv] (130,50) .. controls (138.0,50.0) and (138.0,62.0) .. (130,62.0) .. controls (138.0,62.0) and (138.0,74.0) .. (130,74.0) .. controls (138.0,74.0) and (138.0,86.0) .. (130,86.0) .. controls (138.0,86.0) and (138.0,98.0) .. (130,98.0) .. controls (138.0,98.0) and (138.0,110.0) .. (130,110.0) .. controls (138.0,110.0) and (138.0,122.0) .. (130,122.0) .. controls (138.0,122.0) and (138.0,134.0) .. (130,134.0);
  \draw[lv] (170,50) .. controls (162.0,50.0) and (162.0,62.0) .. (170,62.0) .. controls (162.0,62.0) and (162.0,74.0) .. (170,74.0) .. controls (162.0,74.0) and (162.0,86.0) .. (170,86.0) .. controls (162.0,86.0) and (162.0,98.0) .. (170,98.0) .. controls (162.0,98.0) and (162.0,110.0) .. (170,110.0) .. controls (162.0,110.0) and (162.0,122.0) .. (170,122.0) .. controls (162.0,122.0) and (162.0,134.0) .. (170,134.0);
  \draw[core] (148.0,44.0) -- (148.0,140.0);
  \draw[core] (152.0,44.0) -- (152.0,140.0);
  \draw[hv] (130.0,50.0) -- (80.0,50.0);
  \draw[hv] (130.0,134.0) -- (80.0,134.0);
  \fill[atInfo] (80.0,50.0) circle (3.0);
  \fill[atInfo] (80.0,134.0) circle (3.0);
  \node[font=\footnotesize, text=atInfo, anchor=east, inner sep=0.5pt] at (72.0,50.0) {L1};
  \node[font=\footnotesize, text=atInfo, anchor=east, inner sep=0.5pt] at (72.0,134.0) {L2};
  \draw[aux] (44.0,50.0) -- (44.0,134.0);
  \node[font=\footnotesize, text=atInfo, anchor=center, inner sep=0.5pt, rotate=90] at (34.0,92.0) {400\,V};
  \draw[lv] (170.0,50.0) -- (330.0,50.0);
  \draw[lv] (170.0,134.0) -- (330.0,134.0);
  \fill[atAccentFg] (330.0,50.0) circle (3.0);
  \fill[atAccentFg] (330.0,134.0) circle (3.0);
  \node[font=\footnotesize, text=atAccentFg, anchor=west, inner sep=0.5pt] at (338.0,50.0) {L};
  \node[font=\footnotesize, text=atAccentFg, anchor=west, inner sep=0.5pt] at (338.0,134.0) {N};
  \draw[aux] (370.0,50.0) -- (370.0,134.0);
  \node[font=\footnotesize, text=atAccentFg, anchor=center, inner sep=0.5pt, rotate=-90] at (380.0,92.0) {230\,V};
  \draw[box] (226.0,76.0) rectangle (272.0,108.0);
  \node[font=\footnotesize, text=atFg, anchor=center, inner sep=0.5pt] at (249.0,92.0) {IMD};
  \draw[ph] (240.0,50.0) -- (240.0,76.0);
  \draw[ph] (258.0,134.0) -- (258.0,108.0);
  \draw[pey] (249.0,108.0) -- (249.0,170.0) -- (330.0,170.0);
  \draw[peg] (249.0,108.0) -- (249.0,170.0) -- (330.0,170.0);
  \node[font=\footnotesize, text=atSuccess, anchor=west, inner sep=0.5pt] at (338.0,170.0) {PE};
\end{tikzpicture}

  \caption{Single-phase isolating transformer 400\,V -- 230\,V supplying an
  IT system with insulation monitoring device (IMD).}
  \label{fig:de-it}
\end{figure}
\FloatBarrier

\subsection{Earthing systems}
\label{sec:de-earthing}

DIN~VDE~0100-100 designates the earthing system by two
letters~[20]. The first gives the relation of the source to
earth: T, one point (normally the neutral) directly earthed; I, all live
parts isolated from earth or one point earthed through a high impedance. The
second gives the earthing of the exposed conductive parts of the
installation: N, connected to the earthed point of the source through
protective conductors; T, connected to an earth electrode independent of
the source earth. For TN systems, C denotes a combined protective earth and
neutral conductor (PEN) and S separate N and PE conductors.
Figure~\ref{fig:de-earthing} shows the five systems.

\begin{figure}[!ht]
  \centering
  % generated by north-american-practice/fig/connection_figures.py -- do not edit
\begin{tikzpicture}[x=0.22mm, y=-0.22mm, line cap=round, line join=round]
  \draw[box] (10.0,20.0) rectangle (44.0,104.0);
  \node[font=\scriptsize, text=atFgMuted, anchor=center, inner sep=0.5pt, rotate=90] at (27.0,62.0) {Dyn5};
  \node[font=\small\bfseries, text=atFg, anchor=west, inner sep=0.5pt] at (10.0,6.0) {TN-S};
  \draw[ph] (44.0,30.0) -- (300.0,30.0);
  \node[font=\scriptsize, text=atFgMuted, anchor=west, inner sep=0.5pt] at (305.0,30.0) {L1};
  \draw[ph] (44.0,44.0) -- (300.0,44.0);
  \node[font=\scriptsize, text=atFgMuted, anchor=west, inner sep=0.5pt] at (305.0,44.0) {L2};
  \draw[ph] (44.0,58.0) -- (300.0,58.0);
  \node[font=\scriptsize, text=atFgMuted, anchor=west, inner sep=0.5pt] at (305.0,58.0) {L3};
  \draw[box] (200.0,118.0) rectangle (278.0,152.0);
  \node[font=\scriptsize, text=atFgMuted, anchor=center, inner sep=0.5pt] at (239.0,135.0) {equipment};
  \draw[ph] (214.0,30.0) -- (214.0,118.0);
  \fill[atFg] (214.0,30.0) circle (3.0);
  \draw[nl] (44.0,76.0) -- (300.0,76.0);
  \draw[pey] (56.0,94.0) -- (300.0,94.0);
  \draw[peg] (56.0,94.0) -- (300.0,94.0);
  \draw[pey] (56.0,76.0) -- (56.0,94.0);
  \node[font=\scriptsize, text=atInfo, anchor=west, inner sep=0.5pt] at (305.0,76.0) {N};
  \node[font=\scriptsize, text=atSuccess, anchor=west, inner sep=0.5pt] at (305.0,94.0) {PE};
  \fill[atInfo] (56.0,76.0) circle (3.0);
  \draw[gnd] (56.0,76.0) -- (56.0,140.0);
  \draw[gnd] (56.0,140.0) -- (56.0,150.0);
  \draw[gnd] (49.0,150.0) -- (63.0,150.0);
  \draw[gnd] (51.5,154.0) -- (60.5,154.0);
  \draw[gnd] (54.0,158.0) -- (58.0,158.0);
  \draw[nl] (234.0,76.0) -- (234.0,118.0);
  \fill[atInfo] (234.0,76.0) circle (3.0);
  \draw[pey] (262.0,94.0) -- (262.0,118.0);
  \draw[peg] (262.0,94.0) -- (262.0,118.0);
  \draw[box] (370.0,20.0) rectangle (404.0,104.0);
  \node[font=\scriptsize, text=atFgMuted, anchor=center, inner sep=0.5pt, rotate=90] at (387.0,62.0) {Dyn5};
  \node[font=\small\bfseries, text=atFg, anchor=west, inner sep=0.5pt] at (370.0,6.0) {TN-C};
  \draw[ph] (404.0,30.0) -- (660.0,30.0);
  \node[font=\scriptsize, text=atFgMuted, anchor=west, inner sep=0.5pt] at (665.0,30.0) {L1};
  \draw[ph] (404.0,44.0) -- (660.0,44.0);
  \node[font=\scriptsize, text=atFgMuted, anchor=west, inner sep=0.5pt] at (665.0,44.0) {L2};
  \draw[ph] (404.0,58.0) -- (660.0,58.0);
  \node[font=\scriptsize, text=atFgMuted, anchor=west, inner sep=0.5pt] at (665.0,58.0) {L3};
  \draw[box] (560.0,118.0) rectangle (638.0,152.0);
  \node[font=\scriptsize, text=atFgMuted, anchor=center, inner sep=0.5pt] at (599.0,135.0) {equipment};
  \draw[ph] (574.0,30.0) -- (574.0,118.0);
  \fill[atFg] (574.0,30.0) circle (3.0);
  \draw[nl] (404.0,76.0) -- (660.0,76.0);
  \draw[peg, line width=0.8pt] (404.0,76.0) -- (660.0,76.0);
  \node[font=\scriptsize, text=atSuccess, anchor=west, inner sep=0.5pt] at (665.0,76.0) {PEN};
  \fill[atInfo] (416.0,76.0) circle (3.0);
  \draw[gnd] (416.0,76.0) -- (416.0,140.0);
  \draw[gnd] (416.0,140.0) -- (416.0,150.0);
  \draw[gnd] (409.0,150.0) -- (423.0,150.0);
  \draw[gnd] (411.5,154.0) -- (420.5,154.0);
  \draw[gnd] (414.0,158.0) -- (418.0,158.0);
  \draw[nl] (594.0,76.0) -- (594.0,118.0);
  \fill[atInfo] (594.0,76.0) circle (3.0);
  \draw[peg, line width=0.8pt] (622.0,76.0) -- (622.0,118.0);
  \fill[atInfo] (622.0,76.0) circle (3.0);
  \draw[box] (10.0,220.0) rectangle (44.0,304.0);
  \node[font=\scriptsize, text=atFgMuted, anchor=center, inner sep=0.5pt, rotate=90] at (27.0,262.0) {Dyn5};
  \node[font=\small\bfseries, text=atFg, anchor=west, inner sep=0.5pt] at (10.0,206.0) {TN-C-S};
  \draw[ph] (44.0,230.0) -- (300.0,230.0);
  \node[font=\scriptsize, text=atFgMuted, anchor=west, inner sep=0.5pt] at (305.0,230.0) {L1};
  \draw[ph] (44.0,244.0) -- (300.0,244.0);
  \node[font=\scriptsize, text=atFgMuted, anchor=west, inner sep=0.5pt] at (305.0,244.0) {L2};
  \draw[ph] (44.0,258.0) -- (300.0,258.0);
  \node[font=\scriptsize, text=atFgMuted, anchor=west, inner sep=0.5pt] at (305.0,258.0) {L3};
  \draw[box] (200.0,318.0) rectangle (278.0,352.0);
  \node[font=\scriptsize, text=atFgMuted, anchor=center, inner sep=0.5pt] at (239.0,335.0) {equipment};
  \draw[ph] (214.0,230.0) -- (214.0,318.0);
  \fill[atFg] (214.0,230.0) circle (3.0);
  \draw[nl] (44.0,276.0) -- (300.0,276.0);
  \draw[peg, line width=0.8pt] (44.0,276.0) -- (150.0,276.0);
  \draw[pey] (150.0,294.0) -- (300.0,294.0);
  \draw[peg] (150.0,294.0) -- (300.0,294.0);
  \draw[pey] (150.0,276.0) -- (150.0,294.0);
  \draw[peg] (150.0,276.0) -- (150.0,294.0);
  \fill[atFg] (150.0,276.0) circle (3.0);
  \node[font=\scriptsize, text=atFgMuted, anchor=center, inner sep=0.5pt] at (150.0,268.0) {split};
  \node[font=\scriptsize, text=atSuccess, anchor=center, inner sep=0.5pt] at (86.0,285.0) {PEN};
  \node[font=\scriptsize, text=atInfo, anchor=west, inner sep=0.5pt] at (305.0,276.0) {N};
  \node[font=\scriptsize, text=atSuccess, anchor=west, inner sep=0.5pt] at (305.0,294.0) {PE};
  \fill[atInfo] (56.0,276.0) circle (3.0);
  \draw[gnd] (56.0,276.0) -- (56.0,340.0);
  \draw[gnd] (56.0,340.0) -- (56.0,350.0);
  \draw[gnd] (49.0,350.0) -- (63.0,350.0);
  \draw[gnd] (51.5,354.0) -- (60.5,354.0);
  \draw[gnd] (54.0,358.0) -- (58.0,358.0);
  \draw[nl] (234.0,276.0) -- (234.0,318.0);
  \fill[atInfo] (234.0,276.0) circle (3.0);
  \draw[gnd] (150.0,294.0) -- (150.0,340.0);
  \draw[gnd] (150.0,340.0) -- (150.0,350.0);
  \draw[gnd] (143.0,350.0) -- (157.0,350.0);
  \draw[gnd] (145.5,354.0) -- (154.5,354.0);
  \draw[gnd] (148.0,358.0) -- (152.0,358.0);
  \draw[pey] (262.0,294.0) -- (262.0,318.0);
  \draw[peg] (262.0,294.0) -- (262.0,318.0);
  \draw[box] (370.0,220.0) rectangle (404.0,304.0);
  \node[font=\scriptsize, text=atFgMuted, anchor=center, inner sep=0.5pt, rotate=90] at (387.0,262.0) {Dyn5};
  \node[font=\small\bfseries, text=atFg, anchor=west, inner sep=0.5pt] at (370.0,206.0) {TT};
  \draw[ph] (404.0,230.0) -- (660.0,230.0);
  \node[font=\scriptsize, text=atFgMuted, anchor=west, inner sep=0.5pt] at (665.0,230.0) {L1};
  \draw[ph] (404.0,244.0) -- (660.0,244.0);
  \node[font=\scriptsize, text=atFgMuted, anchor=west, inner sep=0.5pt] at (665.0,244.0) {L2};
  \draw[ph] (404.0,258.0) -- (660.0,258.0);
  \node[font=\scriptsize, text=atFgMuted, anchor=west, inner sep=0.5pt] at (665.0,258.0) {L3};
  \draw[box] (560.0,318.0) rectangle (638.0,352.0);
  \node[font=\scriptsize, text=atFgMuted, anchor=center, inner sep=0.5pt] at (599.0,335.0) {equipment};
  \draw[ph] (574.0,230.0) -- (574.0,318.0);
  \fill[atFg] (574.0,230.0) circle (3.0);
  \draw[nl] (404.0,276.0) -- (660.0,276.0);
  \node[font=\scriptsize, text=atInfo, anchor=west, inner sep=0.5pt] at (665.0,276.0) {N};
  \fill[atInfo] (416.0,276.0) circle (3.0);
  \draw[gnd] (416.0,276.0) -- (416.0,340.0);
  \draw[gnd] (416.0,340.0) -- (416.0,350.0);
  \draw[gnd] (409.0,350.0) -- (423.0,350.0);
  \draw[gnd] (411.5,354.0) -- (420.5,354.0);
  \draw[gnd] (414.0,358.0) -- (418.0,358.0);
  \draw[nl] (594.0,276.0) -- (594.0,318.0);
  \fill[atInfo] (594.0,276.0) circle (3.0);
  \draw[pey] (622.0,352.0) -- (622.0,360.0) -- (660.0,360.0);
  \draw[peg] (622.0,352.0) -- (622.0,360.0) -- (660.0,360.0);
  \draw[gnd] (660.0,360.0) -- (660.0,370.0);
  \draw[gnd] (653.0,370.0) -- (667.0,370.0);
  \draw[gnd] (655.5,374.0) -- (664.5,374.0);
  \draw[gnd] (658.0,378.0) -- (662.0,378.0);
  \node[font=\scriptsize, text=atSuccess, anchor=center, inner sep=0.5pt] at (652.0,350.0) {local};
  \draw[box] (10.0,420.0) rectangle (44.0,504.0);
  \node[font=\scriptsize, text=atFgMuted, anchor=center, inner sep=0.5pt, rotate=90] at (27.0,462.0) {Dyn5};
  \node[font=\small\bfseries, text=atFg, anchor=west, inner sep=0.5pt] at (10.0,406.0) {IT};
  \draw[ph] (44.0,430.0) -- (300.0,430.0);
  \node[font=\scriptsize, text=atFgMuted, anchor=west, inner sep=0.5pt] at (305.0,430.0) {L1};
  \draw[ph] (44.0,444.0) -- (300.0,444.0);
  \node[font=\scriptsize, text=atFgMuted, anchor=west, inner sep=0.5pt] at (305.0,444.0) {L2};
  \draw[ph] (44.0,458.0) -- (300.0,458.0);
  \node[font=\scriptsize, text=atFgMuted, anchor=west, inner sep=0.5pt] at (305.0,458.0) {L3};
  \draw[box] (200.0,518.0) rectangle (278.0,552.0);
  \node[font=\scriptsize, text=atFgMuted, anchor=center, inner sep=0.5pt] at (239.0,535.0) {equipment};
  \draw[ph] (214.0,430.0) -- (214.0,518.0);
  \fill[atFg] (214.0,430.0) circle (3.0);
  \draw[box] (47.0,482.0) rectangle (65.0,504.0);
  \node[font=\scriptsize, text=atFgMuted, anchor=center, inner sep=0.5pt] at (56.0,493.0) {Z};
  \draw[gnd] (56.0,458.0) -- (56.0,482.0);
  \draw[gnd] (56.0,504.0) -- (56.0,540.0);
  \draw[gnd] (56.0,540.0) -- (56.0,550.0);
  \draw[gnd] (49.0,550.0) -- (63.0,550.0);
  \draw[gnd] (51.5,554.0) -- (60.5,554.0);
  \draw[gnd] (54.0,558.0) -- (58.0,558.0);
  \fill[atFg] (56.0,458.0) circle (3.0);
  \node[font=\scriptsize, text=atFgMuted, anchor=west, inner sep=0.5pt] at (70.0,493.0) {isolated or high impedance};
  \draw[ph] (234.0,444.0) -- (234.0,518.0);
  \fill[atFg] (234.0,444.0) circle (3.0);
  \draw[pey] (262.0,552.0) -- (262.0,560.0) -- (300.0,560.0);
  \draw[peg] (262.0,552.0) -- (262.0,560.0) -- (300.0,560.0);
  \draw[gnd] (300.0,560.0) -- (300.0,570.0);
  \draw[gnd] (293.0,570.0) -- (307.0,570.0);
  \draw[gnd] (295.5,574.0) -- (304.5,574.0);
  \draw[gnd] (298.0,578.0) -- (302.0,578.0);
  \node[font=\scriptsize, text=atSuccess, anchor=center, inner sep=0.5pt] at (292.0,550.0) {local};
\end{tikzpicture}

  \caption{Earthing systems: source on the left, equipment with its exposed
  conductive part on the right. Blue: neutral; green-yellow: protective
  earth; green dashed on blue: PEN.}
  \label{fig:de-earthing}
\end{figure}

\subsubsection{TN-S}
N and PE are separate from the transformer onwards and are connected only
at the earthed neutral point of the source. No operating current flows on
the PE, which is the reason TN-S is required or recommended for
installations with extensive information technology
equipment~[21].

\subsubsection{TN-C}
A combined PEN conductor serves as neutral and protective conductor. A PEN
conductor requires a cross-section of at least 10\,mm$^2$ copper or
16\,mm$^2$ aluminium in fixed installations~[22]; inside
buildings TN-C is limited to such conductors.

\subsubsection{TN-C-S}
The PEN conductor of the network is split into N and PE at a point of the
installation, normally the service entrance, where the PE is earthed again;
beyond the split, N and PE are not reconnected. Public low-voltage networks
in Germany are operated as TN-C systems up to the customer's service
entrance and continue as TN-S inside the installation, which makes the
combination TN-C-S~[13].

\subsubsection{TT}
The source neutral is earthed at the transformer, and the exposed
conductive parts are connected to a local earth electrode of the
installation. The earth-fault current returns through the soil, and
protection by automatic disconnection uses residual current devices with
\begin{equation}
  R_{a} \cdot I_{\Delta n} \le 50\,\mathrm{V}
  \label{eq:de-tt}
\end{equation}
where $R_{a}$ is the resistance of the local earth electrode and the
protective conductor, and $I_{\Delta n}$ the rated residual operating
current of the RCD~[23]. Some rural networks are operated as TT
systems.

\subsubsection{IT}
The source is isolated from earth or earthed through a high impedance, and
the exposed conductive parts are earthed locally. A first fault to earth
draws only the capacitive current of the system and does not require
disconnection; it is signalled by an insulation monitoring device. A
second fault on another phase is a short circuit and is disconnected as in
a TN or TT system~[23].
\FloatBarrier

\subsection{Comparison with North American practice}
\label{sec:de-comparison}

Table~\ref{tab:de-comparison} compares the connections of
Section~\ref{sec:na} with German practice.

\begin{table}[!ht]
  \centering
  \small
  \begin{tabularx}{\linewidth}{@{}>{\raggedright\arraybackslash}p{0.22\linewidth}
                                 >{\raggedright\arraybackslash}X
                                 >{\raggedright\arraybackslash}X@{}}
    \toprule
    \multicolumn{1}{c}{\textbf{Item}} & \textbf{North America (480\,V to 240\,V)} & \textbf{Germany} \\
    \midrule
    Utilization voltages   & 480Y/277, 240, 208Y/120\,V; 120/240\,V single-phase & 400/230\,V; 690\,V in industry \\
    Frequency              & 60\,Hz & 50\,Hz \\
    Single-phase supply    & 120/240\,V split-phase, centre-tapped & 230\,V line to neutral \\
    Location of the LV transformer & inside the building, from 480\,V & MV/LV substation; LV-to-LV is the exception \\
    Delta--wye phase shift & lower-voltage side lags 30\,$^{\circ}$ (Dyn1) & Dyn5, 150\,$^{\circ}$ \\
    Terminal markings      & H1--H3, X0--X3 & 1U/1V/1W, 2U/2V/2W, 2N \\
    Unearthed systems      & ungrounded or corner-grounded delta & IT system with insulation monitoring \\
    Centre-tapped delta    & in use & not used \\
    Rules                  & NEC (NFPA~70), IEEE~C57.12 series & DIN~VDE~0100, VDE-AR-N~4100/4110, EN~50588-1, IEC~60076 \\
    \bottomrule
  \end{tabularx}
  \caption{North American and German low-voltage transformer practice.}
  \label{tab:de-comparison}
\end{table}
\FloatBarrier

